\begin{table}[t]\centering\caption{Main results on CartPole (return, max 500; mean $\pm$ std over 5 seeds). \texttt{ret\_ood} is the mean return at shift half-widths 1.5 and 2.}\label{tab:main}\resizebox{\columnwidth}{!}{\begin{tabular}{lccccc}
\toprule
Method & ret\_nominal $\uparrow$ & ret\_w100 $\uparrow$ & ret\_ood $\uparrow$ & ret\_robust $\uparrow$ & $n$ \\
\midrule
No randomisation & \textbf{500.00 $\pm$ 0.00} & 389.12 $\pm$ 28.75 & 243.19 $\pm$ 21.85 & 386.59 $\pm$ 16.59 & 5 \\
Uniform DR narrow (w=0.25) & \underline{500.00 $\pm$ 0.00} & 440.18 $\pm$ 29.22 & 292.36 $\pm$ 11.34 & 418.10 $\pm$ 10.98 & 5 \\
Uniform DR wide (w=1.0) & 500.00 $\pm$ 0.00 & \underline{498.47 $\pm$ 1.71} & \underline{411.98 $\pm$ 34.30} & \underline{470.41 $\pm$ 11.70} & 5 \\
\textbf{Success-gated curriculum DR} & 500.00 $\pm$ 0.00 & \textbf{499.27 $\pm$ 0.75} & \textbf{412.89 $\pm$ 38.93} & \textbf{470.84 $\pm$ 13.07} & 5 \\
\bottomrule
\end{tabular}
}\end{table}

Table~\ref{tab:main} gives the main comparison. All four systems reach the 500-step cap on nominal dynamics, so nominal return does not separate them. Under shifted dynamics, no randomisation is clearly the weakest (\texttt{ret\_robust} 386.6), narrow DR is in between (418.1), and wide DR and the curriculum are the best and indistinguishable (470.4 and 470.8). The curriculum is better than no DR on \texttt{ret\_robust} (paired $p=0.003$) and than narrow DR (paired $p=0.004$) but not different from wide DR (paired $p=0.96$); on \texttt{ret\_ood} the curriculum and wide DR also agree (412.9 vs 412.0, $p=0.97$). With 5 seeds a non-significant difference is not evidence of equivalence, but the gap here is below one return unit in 500. Figure~\ref{fig:profile} shows the return as a function of the evaluation shift: all systems tie at $w=0.25$; wherever they separate ($w\ge0.75$) the ordering is no DR $<$ narrow DR $<$ wide DR $\approx$ curriculum, and the benefit of randomisation grows with the size of the shift.

\begin{figure}[t]\centering\includegraphics[width=0.85\columnwidth]{profile_main.pdf}\caption{Return versus evaluation shift half-width for the four systems (mean $\pm$ std over 5 seeds).}\label{fig:profile}\end{figure}

\textbf{Hypotheses.} H1 (wide DR beats no DR on shifted dynamics) is supported (\texttt{ret\_robust} 470.4 vs 386.6, paired $p=0.0007$ from \texttt{rh compare} with wide DR as reference). H2 (a width beyond which nominal return falls while robustness stops improving) is refuted in this setting: see Section~\ref{sec:ablations}. H3 (curriculum matches wide DR) is supported in the sense of no detectable difference; the curriculum does not improve on wide DR, and it also has no cost in robustness.

\textbf{What the numbers do not show.} The easy cap (500 steps) means nominal return cannot rise above the other systems and may hide small nominal costs of wide ranges that a harder objective (e.g.\ tracking error or energy) would reveal. Because the widest evaluation level is $w=2$, systems trained at $w\ge 1.5$ are largely in-range there; true out-of-range behaviour of wide-trained policies is not measured.
